\documentclass[11pt,a4paper]{article}

\usepackage[utf8]{inputenc}
\usepackage[T1]{fontenc}
\usepackage[margin=1in]{geometry}
\usepackage{amsmath,amssymb}
\usepackage{booktabs}
\usepackage{array}
\usepackage{tabularx}
\usepackage{xcolor}
\usepackage{enumitem}
\usepackage{hyperref}
\usepackage{caption}
\usepackage{graphicx}
\usepackage{adjustbox}
\usepackage[final,protrusion=true,expansion=false]{microtype}
\usepackage{float}
\usepackage{multirow}
\usepackage{longtable}

\hypersetup{colorlinks=true,linkcolor=blue!60!black,citecolor=blue!60!black,urlcolor=blue!60!black}
\setlength{\emergencystretch}{2em}
\setlist[itemize]{leftmargin=1.2em,itemsep=1pt,topsep=2pt}
\setlist[enumerate]{leftmargin=1.2em,itemsep=1pt,topsep=2pt}

\newcommand{\up}{\ensuremath{\uparrow}}
\newcommand{\down}{\ensuremath{\downarrow}}

\title{\textbf{TF-OVOS}: A Four-Part Evaluation Framework for \\
       Training-Free Open-Vocabulary Object Segmentation \\
       \vspace{0.3em}
       \large{Proposal Draft v25 -- two-tier metrics with expanded exploratory set}}
\author{Draft for Discussion}
\date{\today}

\begin{document}
\maketitle

%==============================================================
\section{Core Idea}

This proposal studies \textbf{training-free open-vocabulary object segmentation} (TF-OVOS). Given an image $I$ and a candidate vocabulary $\mathcal{C}$, a method predicts a set of mask--category pairs with optional confidence scores:
\begin{equation}
I+\mathcal{C}\rightarrow \mathcal{P}=\{(\hat{M}_k,\hat{y}_k,\hat{s}_k)\}_{k=1}^{K}, \qquad \hat{y}_k\in\mathcal{C}.
\end{equation}
For semantic-segmentation benchmarks, $\mathcal{P}$ is converted into a dense class map. For object-centric or hard-domain targets, the same output is converted into one or more foreground object masks. This makes OVCamo and other camouflaged-object datasets \emph{hard-domain targets} rather than the whole topic.

\paragraph{Strict training-free setting.}
The main rows are \textbf{training-free and MLLM/API-free}: no target-dataset training, no prompt tuning, no adapter training, no validation threshold tuning, no inference-time generative MLLM calls, and no closed-source API-VLMs. Frozen public models such as CLIP, SigLIP, GroundingDINO, DINOv2, SAM, and SAM2 are allowed.

\paragraph{Compact trained references.}
Training-based rows are kept intentionally small. The main paper keeps only a few general open-vocabulary segmentation references, such as OVSeg~\cite{ovseg2023}, SAN~\cite{san2023}, and ODISE~\cite{odise2023}. They are not strict-TF competitors; they provide context for effectiveness, source-to-target generalization, and efficiency.

%==============================================================

\section{Evaluation Protocol}

Following advisor feedback, the revised framework adopts a \textbf{two-tier metric design}.
The \emph{standard tier} keeps \textbf{one common standard metric family} (mIoU) for E1--E3 whenever semantic labels are available, so the leaderboard, robustness, and cross-dataset comparisons all use the same primary number; the evaluations differ mainly by \emph{dataset difficulty and scenario}, not by inventing a new ranking metric for each block. The \emph{exploratory tier} records a broader set of readouts that are not used for ranking but are explicitly intended to surface where standard mIoU is silent --- naming vs.\ localization failures, ambiguity in large vocabularies, proposal recall, mask--category consistency, and stuff/thing skew. The exploratory tier is what the proposal expects to drive novelty: any new finding or method contribution must first appear as a pattern in these readouts before it is promoted into the main paper. The two tiers are reported separately so that the standard tier stays clean and comparable while the exploratory tier remains a discovery tool rather than a moving ranking target.

\begin{table}[!htbp]
\centering
\small
\caption{Final E1--E4 evaluation protocol. E1--E3 share mIoU as the common \emph{standard} primary metric whenever semantic labels are available. The \emph{exploratory} column records additional readouts intended to find phenomena that mIoU alone cannot reveal; these are not used to rank methods. E4 is separated because it measures cost rather than segmentation quality.}
\label{tab:evaluation-overview}
\setlength{\tabcolsep}{3pt}
\begin{adjustbox}{max width=\textwidth}
\begin{tabular}{p{0.13\textwidth} p{0.24\textwidth} p{0.17\textwidth} p{0.30\textwidth} p{0.13\textwidth}}
\toprule
Evaluation & Scenario / dataset difficulty & Standard metric (ranking) & Exploratory readouts (discovery, no ranking) & Role \\
\midrule
\textbf{E1 Standard OVOS effectiveness} &
Mainstream open-vocabulary segmentation datasets with compact official vocabularies: PASCAL VOC-20, PASCAL Context-59, ADE20K-150, and COCO-Stuff-171. &
mIoU on each dataset and average mIoU. &
Per-class IoU, thing/stuff split, pixel accuracy, mean class accuracy, prediction coverage, frequency-stratified mIoU, head/mid/tail-frequency mIoU, small/medium/large object mIoU, confidence calibration (ECE / Brier), empty-prediction and over-segmentation rates. &
Primary effectiveness leaderboard. \\
\addlinespace[2pt]
\textbf{E2 Vocabulary / category-ambiguity robustness} &
Harder large-vocabulary versions of standard datasets: PASCAL Context-459 and ADE20K-847, paired with Context-59 and ADE20K-150. &
mIoU on compact and large vocabularies. &
Proposal-defined diagnostic summaries $\Delta_{\mathrm{vocab}}$, MCMR@0.5/0.75, semantic-distance-weighted mismatch, hierarchical IoU, synonym-collapsed mIoU; prompt-sensitivity variance; vocabulary-size scaling curves (subsampled sizes); region-level Top-1/Top-3/Top-5 naming accuracy; top mismatch pairs, confusion communities, novel-vs-base class split; GT-region naming and GT-text localization probes. &
Robustness to vocabulary difficulty. \\
\addlinespace[2pt]
\textbf{E3 Cross-dataset generalization} &
Source-to-target comparison across the same standard semantic targets. Strict-TF methods use no source training; trained references use their official source training. &
Target mIoU on VOC-20, Context-59, ADE20K-150, and COCO-Stuff-171. &
Rank stability across targets, Kendall's $\tau$ / Spearman rank correlation between E1 and E3 rankings, worst-target mIoU, worst-class mIoU within target, source/target gap, per-target $\Delta$ vs.\ E1, per-target $\Delta$ MCMR, class-overlap-stratified mIoU (source-overlapping vs.\ source-novel classes), domain-shift sensitivity (mIoU drop vs.\ embedding-distance to source), transferability of MCMR@0.5. &
Generalization comparison. \\
\addlinespace[2pt]
\textbf{E4 Efficiency comparison} &
Same method rows under fixed hardware, input size, and adapter settings. &
Not applicable. &
Training time, offline cost, inference time, peak memory, trainable parameters, model calls, cost/mIoU trade-off, mIoU per GFLOP, mIoU per second, offline cost amortization curve, per-stage memory and latency breakdown. &
Cost analysis. \\
\bottomrule
\end{tabular}
\end{adjustbox}
\end{table}

\paragraph{Two-tier metric rule.}
E1, E2, and E3 all report mIoU as the primary quantitative metric on semantic targets, so the leaderboard remains comparable across scenarios. E2 becomes harder because the vocabulary is larger and more ambiguous; E3 becomes harder because the evaluation emphasizes source-to-target transfer. In parallel, a broader exploratory metric set is recorded on every run (see Section~\ref{sec:metrics}); these readouts do not change the ranking but are the channel through which the proposal expects to identify failure modes, justify novelty, and decide which patterns are worth a dedicated section in the paper. An exploratory readout is promoted into the main paper only after it reveals a stable, method-discriminating pattern across multiple datasets.

\paragraph{Why this replaces the previous OVCOS-centered structure.}
The previous draft was too narrow because OVCamo was both the task definition and the main target. The revised E1--E4 framework makes OVCamo one optional hard-domain transfer target, while the main topic becomes general training-free open-vocabulary object / semantic segmentation.

\begin{table}[!htbp]
\centering
\scriptsize
\caption{Main dataset plan. E1--E3 use standard semantic-segmentation targets so that the primary metric can remain mIoU. Camouflaged-object datasets are kept as appendix hard-domain stress targets because their metrics are not fully aligned with standard OVSS mIoU.}
\label{tab:datasets}
\setlength{\tabcolsep}{3pt}
\begin{adjustbox}{max width=\textwidth}
\begin{tabular}{p{0.18\textwidth} p{0.22\textwidth} p{0.31\textwidth} p{0.20\textwidth}}
\toprule
Dataset / split & Role & What it provides & Primary metric \\
\midrule
PASCAL VOC-20~\cite{pascalvoc2010} & E1 and E3 standard target & Object-centric semantic masks and 20 foreground classes. & mIoU. \\
PASCAL Context-59~\cite{pascalcontext2014} & E1, E2 compact-vocabulary side, and E3 target & Scene semantic masks with object and stuff categories. & mIoU. \\
PASCAL Context-459~\cite{pascalcontext2014} & E2 large-vocabulary side & Larger label space for vocabulary robustness and category-ambiguity analysis. & mIoU. \\
ADE20K-150~\cite{ade20k2017} & E1, E2 compact-vocabulary side, and E3 target & Diverse scene parsing categories. & mIoU. \\
ADE20K-847~\cite{ade20k2017} & E2 large-vocabulary side & Larger open-vocabulary label space. & mIoU. \\
COCO-Stuff-171~\cite{cocostuff2018} & E1 and E3 standard target; common source for trained refs & Object and stuff semantic categories; frequently used by trained OVSS methods. & mIoU. \\
\midrule
OVCamo-TE~\cite{ovcoser2024}, CAMO-TE~\cite{camo2019}, COD10K-TE-Camo~\cite{cod10k2020}, NC4K~\cite{nc4k2021} & Appendix hard-domain transfer targets & Camouflaged-object masks; only OVCamo has aligned OVCOS class labels. & Mask-only metrics; OVCamo c-metrics optional. \\
\bottomrule
\end{tabular}
\end{adjustbox}
\end{table}

\paragraph{Hard-domain target rule.}
Hard-domain camouflaged-object datasets are useful for stress testing, but they are moved out of the core E1--E3 metric-unified tables because CAMO, COD10K, and NC4K do not provide the same open-vocabulary semantic label space. They are reported as appendix transfer targets with mask-only metrics. This keeps the main benchmark clean while preserving a place to test whether standard OVOS performance transfers to difficult object domains.

%==============================================================
\section{Metrics}
\label{sec:metrics}

\paragraph{Two-tier philosophy.}
The benchmark explicitly distinguishes two tiers of metrics. The \textbf{standard tier} is a single, common mIoU family used for E1--E3 ranking and reporting; it is what makes the leaderboard comparable. The \textbf{exploratory tier} is a broader set of readouts recorded on every run; it is what the proposal expects to use to find phenomena that mIoU cannot expose by itself. The two tiers serve different purposes and must not be conflated: the standard tier protects comparability, while the exploratory tier protects novelty. Concretely, the plan is to first run the standard metric on every method, and \emph{in the same evaluation pass} record the exploratory readouts, so any interesting pattern (a method that ranks well on mIoU but breaks under MCMR; a method whose mIoU is stable but whose stuff/thing split collapses; a method whose proposal-recall ceiling explains its mIoU floor) becomes available without re-running the pipeline.

\paragraph{Common primary (standard) metric for E1--E3.}
For all standard semantic-segmentation targets, the primary metric is mIoU. The same metric is used in E1, E2, and E3; the evaluations differ only in dataset/vocabulary difficulty and transfer scenario. A benchmark should not need a different main metric for every scenario.

\paragraph{Exploratory metric set.}
The exploratory tier is divided into three groups. (i) \emph{Standard supporting readouts} are well-known metrics from the segmentation literature that are not proposal-defined --- per-class IoU, thing/stuff mIoU, pixel accuracy, mean class accuracy, frequency-stratified mIoU (including head/mid/tail-frequency mIoU), small/medium/large object mIoU, BIoU, proposal Recall@\{0.5, 0.7\}, confidence calibration (ECE / Brier), empty-prediction and over-segmentation rates. They cost almost nothing to compute once mIoU is computed. (ii) \emph{Proposal-defined diagnostic summaries} are introduced by this work to surface vocabulary-driven and mask--category-driven failure modes that the standard tier averages away. The current set is $\Delta_{\mathrm{vocab}}$, MCMR@$\tau$, semantic-distance-weighted mismatch, hierarchical IoU, synonym-collapsed mIoU, prompt-sensitivity variance, vocabulary-size scaling curves, Kendall's $\tau$ rank correlation between E1 and E3, class-overlap-stratified mIoU, and domain-shift sensitivity. (iii) \emph{Probe-style measurements} isolate naming vs.\ localization failure with GT-region naming Top-1/Top-5, region-level Top-$k$ naming accuracy, GT-text localization IoU/BIoU, proposal-recall oracles, and MCC consistency before/after. Groups (ii) and (iii) are the channel through which novelty is expected to surface. None of the exploratory metrics is used to re-rank the leaderboard. They are reported in dedicated columns or appendix tables clearly separated from the standard-tier columns.

\paragraph{Metric status.}
Table~\ref{tab:metric-status} lists every metric category with its tier, status, and intended use. The reader should be able to tell at a glance, for any number reported in the paper, whether it is a ranking number or an exploratory number.

{\scriptsize
\setlength{\tabcolsep}{3pt}
\begin{longtable}{p{0.18\textwidth} p{0.30\textwidth} p{0.16\textwidth} p{0.27\textwidth}}
\caption{Metric status and intended use. The benchmark uses a two-tier design: a single common standard metric for ranking E1--E3, and a broader exploratory set for finding phenomena and supporting novelty. Exploratory metrics never replace mIoU on the leaderboard. Rows are grouped by the evaluation block they primarily inform.}
\label{tab:metric-status}\\
\toprule
Metric & Examples & Tier / status & Intended use \\
\midrule
\endfirsthead
\multicolumn{4}{l}{\emph{Table~\ref{tab:metric-status} continued from previous page.}}\\
\toprule
Metric & Examples & Tier / status & Intended use \\
\midrule
\endhead
\midrule
\multicolumn{4}{r}{\emph{Continued on next page.}}\\
\endfoot
\bottomrule
\endlastfoot
\multicolumn{4}{l}{\textbf{Standard tier (used for E1/E2/E3 ranking)}} \\
Common standard metric & mIoU on VOC-20, Context-59/459, ADE20K-150/847, COCO-Stuff-171 & \textbf{Standard}; primary reporting & Unified effectiveness, robustness, and generalization ranking. \\
\midrule
\multicolumn{4}{l}{\textbf{Exploratory tier --- E1 (effectiveness diagnostics)}} \\
Per-class and split readouts & Per-class IoU, thing/stuff mIoU, pixel accuracy, mean class accuracy & Group (i); cheap & Inspect where mIoU changes come from. \\
Frequency-stratified mIoU & Head/mid/tail-frequency class mIoU (tertile split by training-set frequency) & Group (i); cheap & Detect head-class-only methods; expose long-tail behavior. \\
Object-scale stratified mIoU & Small / medium / large object mIoU (COCO-style scale bins) & Group (i); cheap & Reveal resolution and proposal-granularity bottlenecks. \\
Confidence calibration & ECE, Brier score on confidence vs.\ pixel/region correctness & Group (i); cheap & TF methods rarely report calibration; potential novelty if systematically miscalibrated. \\
Coverage / over-segmentation & Empty-prediction rate, over-segmentation rate, prediction coverage & Group (i); cheap & Detect silent failures and proposal-flooding behavior. \\
\midrule
\multicolumn{4}{l}{\textbf{Exploratory tier --- E2 (vocabulary / ambiguity diagnostics, primary novelty channel)}} \\
$\Delta_{\mathrm{vocab}}$ & mIoU drop from compact $\to$ large official vocabulary (per-dataset and averaged) & \textbf{Group (ii)}; novelty-driving & Summarize vocabulary-robustness gap. \\
MCMR@$\tau$ & MCMR@0.5 (primary), MCMR@0.75 (strict) & \textbf{Group (ii)}; novelty-driving & Localized-region category-mismatch rate; separates naming failure from localization failure. \\
Semantic-distance-weighted mismatch & Mismatch weighted by CLIP/WordNet distance between $\hat{y}_i$ and $y_i$ & \textbf{Group (ii)}; novelty-driving & Distinguish ``near-miss'' (sofa $\to$ couch) from ``wild-miss'' (sofa $\to$ building). \\
Hierarchical IoU & Partial credit when predicted label is a WordNet hypernym/hyponym of GT label & \textbf{Group (ii)}; novelty-driving & Test whether plain mIoU underestimates OV semantic ability. \\
Synonym-collapsed mIoU & mIoU after collapsing synonym classes into single nodes & \textbf{Group (ii)}; novelty-driving & Quantify how much vocabulary-growth drop is real ambiguity vs.\ synonym competition. \\
Prompt-sensitivity variance & mIoU std.\ across a fixed prompt-template set per method & \textbf{Group (ii)}; novelty-driving & Distinguish prompt-fragile CLIP-based methods from prompt-stable ones. \\
Vocabulary-size scaling curve & mIoU at vocabulary sizes \{59, 150, 300, 459\} (Context) and \{150, 300, 500, 847\} (ADE) by random subsampling & \textbf{Group (ii)}; novelty-driving & Curve shape itself is a finding (linear / saturating / cliff). \\
Region-level Top-$k$ naming & Top-1/Top-3/Top-5 accuracy on best-matched regions ($\mathrm{IoU}\ge 0.5$) & \textbf{Group (iii)}; probe & Separate ``does not know'' from ``knows but ranks wrong.'' \\
GT-region naming probe & Top-1/Top-5 naming accuracy on GT-cropped regions & \textbf{Group (iii)}; probe & Isolate naming ability with localization removed. \\
GT-text localization probe & IoU / BIoU when prompted with GT class name & \textbf{Group (iii)}; probe & Isolate localization ability with naming removed. \\
Proposal-recall oracle & BestIoU and Recall@\{0.5, 0.7\} of class-agnostic proposals & \textbf{Group (iii)}; probe & Ceiling for proposal-based pipelines. \\
MCC consistency probe & mIoU / MCMR before vs.\ after post-hoc mask--category re-ranking & \textbf{Group (iii)}; probe & Test whether consistency re-ranking helps. \\
Qualitative ambiguity evidence & Top mismatch pairs, confusion communities, novel-vs-base class split & Qualitative & Illustrate systematic label confusions. \\
\midrule
\multicolumn{4}{l}{\textbf{Exploratory tier --- E3 (cross-dataset generalization diagnostics)}} \\
Rank correlations & Kendall's $\tau$ and Spearman $\rho$ between E1 and E3 method rankings; cross-target rank stability & \textbf{Group (ii)}; novelty-driving & Hard evidence of generalization fragility. \\
Class-overlap-stratified mIoU & Target-class mIoU split into source-overlapping vs.\ source-novel classes & \textbf{Group (ii)}; novelty-driving & Directly exposes trained-reference dependence on source vocabulary; favors TF rows on novel classes. \\
Domain-shift sensitivity & mIoU drop vs.\ DINO/CLIP-embedding distance from source to target distribution & \textbf{Group (ii)}; novelty-driving & TF methods may show drop uncorrelated with distance --- a clean TF-vs-trained signature. \\
Worst-class within target & Lowest per-class IoU within each target & Group (i) & Catches ``average is fine but one class collapsed'' patterns. \\
Per-target $\Delta$ MCMR & MCMR@0.5 transferred from E1/E2 datasets to E3 targets & Group (ii) & Separate ``localization broke during transfer'' from ``naming broke.'' \\
Source/target gap & Side-by-side source vs.\ target mIoU for trained refs & Group (i) & Contextualize trained-method generalization claims. \\
\midrule
\multicolumn{4}{l}{\textbf{E4 (efficiency)}} \\
Core efficiency & Training time, offline preprocessing, inference time, peak memory, trainable parameters, model calls & Standard cost reporting & Explain cost behind performance. \\
Normalized efficiency & mIoU per GFLOP, mIoU per second & Group (ii); novelty-supporting & Single normalized number for Pareto-frontier plots. \\
Offline amortization curve & Per-image total cost as a function of evaluation set size $N$ & Group (ii); novelty-supporting & Fair treatment of diffusion-prototype methods (FreeDA, OVDiff). \\
Per-stage breakdown & Memory and latency split across encoder / proposal / naming / post-processing stages & Group (i) & Identify the bottleneck stage per method. \\
\midrule
\multicolumn{4}{l}{\textbf{Appendix}} \\
Hard-domain metrics & IoU, $S_m$, $F^\omega_\beta$, $E_m$, MAE; OVCamo cIoU / c-metrics & Appendix / stress test & Camouflaged-object transfer where standard semantic mIoU is not aligned. \\
\end{longtable}
}

\paragraph{Promotion rule.}
An exploratory metric is promoted into the main paper only if it satisfies all three of the following: (a) it produces a stable ranking across at least two of \{VOC-20, Context-59/459, ADE-150/847, COCO-Stuff-171\}; (b) it discriminates between method families rather than only between random seeds; (c) it surfaces a failure mode that standard mIoU does not already make obvious. Exploratory metrics that fail this rule stay in the spreadsheet and appendix. This is what keeps the exploratory tier honest: it is allowed to be broad precisely because promotion is gated.

\paragraph{E1 standard and exploratory metrics.}
E1 reports mIoU on VOC-20, Context-59, ADE20K-150, COCO-Stuff-171, and their average as the standard ranking number. In parallel, a set of cheap exploratory readouts is recorded so that E1 can answer not just \emph{which method ranks higher} but \emph{on what kinds of inputs the gap actually lives}. These include per-class IoU, pixel accuracy, mean class accuracy, thing/stuff mIoU, and \emph{frequency-stratified mIoU} (each dataset's classes are split into head/mid/tail tertiles by training-set frequency, and mIoU is reported on each tertile). \emph{Object-scale stratified mIoU} reports small/medium/large mIoU using COCO-style scale bins, surfacing resolution and proposal-granularity issues. \emph{Confidence calibration} is recorded as ECE and Brier score on each method's natural confidence output (where defined), since training-free methods rarely report calibration and systematic miscalibration is itself a candidate finding. \emph{Empty-prediction rate}, \emph{over-segmentation rate}, and \emph{prediction coverage} flag silent failure modes that average mIoU hides.

\paragraph{E2 standard and exploratory metrics.}
E2 first reports the same standard metric, mIoU, on compact and large vocabularies. In parallel, it records the two proposal-defined diagnostics.
First, \textbf{vocabulary robustness drop} summarizes how much mIoU degrades when the label space expands from a compact official vocabulary to a larger official vocabulary:
\begin{align}
\Delta_{\mathrm{vocab}}^{\mathrm{Context}} &= \mathrm{mIoU}(\mathrm{Context}{\text -}59)-\mathrm{mIoU}(\mathrm{Context}{\text -}459),\\
\Delta_{\mathrm{vocab}}^{\mathrm{ADE}} &= \mathrm{mIoU}(\mathrm{ADE}{\text -}150)-\mathrm{mIoU}(\mathrm{ADE}{\text -}847),\\
\Delta_{\mathrm{vocab}} &= \frac{1}{2}\left(\Delta_{\mathrm{vocab}}^{\mathrm{Context}}+\Delta_{\mathrm{vocab}}^{\mathrm{ADE}}\right).
\end{align}
A smaller $\Delta_{\mathrm{vocab}}$ means the method is more robust to larger and more ambiguous vocabularies. This summary uses existing dataset splits and standard mIoU; it does not create new labels.

Second, \textbf{mask--category mismatch rate} measures how often a method localizes a region but assigns the wrong category name. For each ground-truth region $G_i$, let $\hat{M}_i$ be the best-matched predicted region and $\hat{y}_i$ its label. Then
\begin{equation}
\mathrm{MCMR@}\tau=
\frac{\sum_i \mathbf{1}[\mathrm{IoU}(\hat{M}_i,G_i)\ge \tau] \mathbf{1}[\hat{y}_i\ne y_i]}
{\max\left(\sum_i \mathbf{1}[\mathrm{IoU}(\hat{M}_i,G_i)\ge \tau],\epsilon\right)}.
\end{equation}
The primary threshold is $\tau=0.5$; $\tau=0.75$ is recorded as a stricter exploratory variant. A lower MCMR means fewer category mismatches after successful localization. Frequent mismatch pairs $y_i\rightarrow\hat{y}_i$ and confusion communities are reported only as qualitative evidence.

E2 additionally records several novelty-oriented exploratory readouts. \emph{Semantic-distance-weighted mismatch} reweights each MCMR error by the CLIP/WordNet distance between $\hat{y}_i$ and $y_i$, separating near-misses (e.g.\ \emph{sofa $\to$ couch}) from wild misses (e.g.\ \emph{sofa $\to$ building}). \emph{Hierarchical IoU} grants partial credit when the predicted label is a WordNet hypernym or hyponym of the GT label; it tests whether plain mIoU systematically underestimates open-vocabulary semantic ability. \emph{Synonym-collapsed mIoU} merges synonym classes (using WordNet synsets or a lightweight manual list on Context-459 / ADE-847) and recomputes mIoU; this isolates how much of the vocabulary-growth drop is genuine ambiguity versus synonym competition over a single visual concept. \emph{Prompt-sensitivity variance} runs each method with a fixed bank of prompt templates (e.g.\ ``a photo of \{X\}'', ``\{X\}'', ``a \{X\} in the scene'') and reports the std.\ of mIoU across templates --- discriminating prompt-fragile CLIP-based methods from prompt-stable detector- or VFM-based ones. \emph{Vocabulary-size scaling curves} subsample the large vocabularies at sizes \{59, 150, 300, 459\} on Context and \{150, 300, 500, 847\} on ADE, plotting mIoU against vocabulary size; the curve's shape (linear, saturating, cliff) is itself the finding. \emph{Region-level Top-$k$ naming accuracy} records, on each well-matched region ($\mathrm{IoU}\ge 0.5$), whether the GT label appears in the predicted Top-1, Top-3, or Top-5 names, separating ``does not know'' from ``knows but ranks wrong.''

\paragraph{E3 standard and exploratory readouts.}
E3 reports target mIoU on standard semantic targets, using the same metric as E1 and E2. To make the generalization behavior visible without adding another ranking metric, the table also records source/training data, target average, worst-target mIoU, \emph{worst-class mIoU within target} (the minimum per-class IoU within each target, catching "average is fine but one class collapsed" patterns), per-target $\Delta$ relative to E1, and the transferability of MCMR@0.5 from compact to larger vocabularies via \emph{per-target $\Delta$ MCMR}. Source and target scores are shown side by side when trained references report source-domain performance.

E3 additionally records three novelty-oriented exploratory readouts. \emph{Rank correlations} compute Kendall's $\tau$ and Spearman's $\rho$ between the E1 method ranking and the E3 per-target rankings; low correlation is hard evidence of generalization fragility that mIoU averages hide. \emph{Class-overlap-stratified mIoU} splits each target's classes into a source-overlapping subset (classes present in or close to the trained reference's source vocabulary) and a source-novel subset, and reports mIoU on each; this directly exposes the trained-reference dependence on source vocabulary and provides a natural slot for strict-TF methods to show their advantage. \emph{Domain-shift sensitivity} computes a DINOv2 or CLIP image-embedding distance between each target dataset and the trained-reference source, then plots target-mIoU drop against this distance; strict-TF rows that show drop uncorrelated with distance produce a clean training-free signature.

\paragraph{E4 cost metrics.}
E4 reports training time, offline preprocessing time, inference time per image, peak GPU memory, trainable parameters, model calls, and module requirements as the core cost numbers. Offline prototype or diffusion stages are reported separately from per-image inference. Three exploratory readouts make the cost analysis more informative. \emph{Normalized efficiency} reports mIoU per GFLOP and mIoU per second, giving a single number for Pareto-frontier plots over the standard mIoU. \emph{Offline cost amortization curves} plot per-image total cost (offline + online) as a function of evaluation-set size $N$; for diffusion-prototype methods (FreeDA, OVDiff) this gives a fair picture where small $N$ punishes them and large $N$ amortizes the offline stage. \emph{Per-stage memory and latency breakdown} splits each method's per-image cost across encoder / proposal / naming / post-processing stages, identifying the bottleneck stage and where future optimization would pay off most.

\section{Methods and Adapter}

The benchmark uses explicit method rows rather than only family names. Final inclusion still requires a code/weight check, but the proposal fixes the intended roster below.

\begin{table}[!htbp]
\centering
\scriptsize
\caption{Concrete method roster and OVOS adapters. Rows marked backup are included if public code/weights are stable enough; otherwise they move to appendix with an implementation note.}
\label{tab:families}
\setlength{\tabcolsep}{3pt}
\begin{adjustbox}{max width=\textwidth}
\begin{tabular}{p{0.17\textwidth} p{0.40\textwidth} p{0.22\textwidth} p{0.17\textwidth}}
\toprule
Family & E1/E2 rows & Default E3 rows & Unified OVOS adapter \\
\midrule
CLIP-only / attention-edited OVSS &
MaskCLIP~\cite{maskclip2022}, CLIP-DIY~\cite{clipdiy2024}, SCLIP~\cite{sclip2024}, NACLIP~\cite{naclip2025}; backup: CLIPtrase~\cite{cliptrase2024}, SC-CLIP~\cite{scclip2025}, ResCLIP~\cite{resclip2025} &
MaskCLIP, SCLIP, NACLIP, ResCLIP &
Convert class score maps into dense predictions or connected-component masks with fixed thresholds. \\
\addlinespace[2pt]
CLIP + VFM OVSS &
ProxyCLIP~\cite{proxyclip2024}, CorrCLIP~\cite{corrclip2025}, Trident~\cite{trident2024}; backup: CASS~\cite{cass2025} &
ProxyCLIP, CorrCLIP, Trident &
Use VFM-refined patch correlations or SAM/DINO spatial priors before dense-map or mask selection. \\
\addlinespace[2pt]
Diffusion / reference-based OVSS &
OVDiff~\cite{ovdiff2024}, FreeDA~\cite{freeda2024} &
FreeDA; OVDiff if code/runtime is feasible &
Use diffusion-generated or diffusion-augmented visual prototypes to match regions with text classes. \\
\addlinespace[2pt]
OV detection + SAM &
GroundingDINO~\cite{groundingdino2023} + SAM~\cite{sam2023}; GroundingDINO + SAM2~\cite{sam22024} &
GroundingDINO + SAM2 &
Prompt the detector with class names, convert boxes into masks with SAM/SAM2, and resolve overlaps by score. \\
\addlinespace[2pt]
Class-agnostic proposal + VLM naming &
SAM-AMG + CLIP/SigLIP~\cite{clip2021,siglip2023}; DINOv2~\cite{dinov22023} + SAM + CLIP/SigLIP; MCC consistency diagnostic &
SAM-AMG + SigLIP; DINOv2 + SAM + SigLIP &
Generate proposals and name proposals with a frozen VLM. MCC is evaluated only as a post-hoc diagnostic probe. \\
\addlinespace[2pt]
Compact trained references &
OVSeg~\cite{ovseg2023}, SAN~\cite{san2023}, ODISE~\cite{odise2023} &
OVSeg, SAN, ODISE &
Not strict-TF. Used as compact context for E1/E3/E4, not as primary benchmark rows. \\
\bottomrule
\end{tabular}
\end{adjustbox}
\end{table}

\paragraph{Adapter rule.}
Every method is converted to either a dense class map or a set of $(\hat{M},\hat{y},\hat{s})$ pairs using pre-declared thresholds, prompts, resize sizes, SAM settings, and post-processing. Constants are fixed before target evaluation and are not tuned per method or per dataset. MCC re-ranking is reported only as a post-hoc consistency diagnostic rather than as a primary leaderboard method or a new method family.

%==============================================================

\section{Main Tables}

\subsection{E1: Standard OVOS Effectiveness}

\begin{table}[H]
\centering
\scriptsize
\caption{E1 standard open-vocabulary object/semantic segmentation effectiveness. Strict-TF rows use no target training. The standard metric is mIoU across all columns. Numeric entries are filled after running the fixed adapters.}
\label{tab:e1-effectiveness}
\setlength{\tabcolsep}{3pt}
\begin{adjustbox}{max width=\textwidth}
\begin{tabular}{l c c c c c c c}
\toprule
Method & Type & Strict TF? & VOC-20 mIoU\up & Context-59 mIoU\up & ADE-150 mIoU\up & COCO-Stuff mIoU\up & Avg.\up \\
\midrule
\multicolumn{8}{l}{\textbf{CLIP-only / attention-edited OVSS}} \\
MaskCLIP & CM & Yes & & & & & \\
CLIP-DIY & CM & Yes & & & & & \\
SCLIP & CM & Yes & & & & & \\
NACLIP & CM & Yes & & & & & \\
CLIPtrase & CM & Yes & & & & & \\
SC-CLIP & CM & Yes & & & & & \\
ResCLIP & CM & Yes & & & & & \\
\midrule
\multicolumn{8}{l}{\textbf{CLIP + VFM OVSS}} \\
ProxyCLIP & VM & Yes & & & & & \\
CorrCLIP & VM & Yes & & & & & \\
Trident & VM & Yes & & & & & \\
CASS & VM & Yes & & & & & \\
\midrule
\multicolumn{8}{l}{\textbf{Diffusion / reference-based OVSS}} \\
OVDiff & DM & Yes & & & & & \\
FreeDA & DM & Yes & & & & & \\
\midrule
\multicolumn{8}{l}{\textbf{OV detection + SAM}} \\
GroundingDINO + SAM & DS & Yes & & & & & \\
GroundingDINO + SAM2 & DS & Yes & & & & & \\
\midrule
\multicolumn{8}{l}{\textbf{Class-agnostic proposal + VLM naming}} \\
SAM-AMG + CLIP & PN & Yes & & & & & \\
SAM-AMG + SigLIP & PN & Yes & & & & & \\
DINOv2 + SAM + CLIP & PN & Yes & & & & & \\
DINOv2 + SAM + SigLIP & PN & Yes & & & & & \\
\midrule
\multicolumn{8}{l}{\textbf{Compact training-based references (not strict TF)}} \\
OVSeg & Trained & No & & & & & \\
SAN & Trained & No & & & & & \\
ODISE & Trained & No & & & & & \\
\bottomrule
\end{tabular}
\end{adjustbox}
\end{table}

\subsection{E2: Vocabulary and Category-Ambiguity Robustness}

\begin{table}[H]
\centering
\scriptsize
\caption{E2 vocabulary and category-ambiguity robustness. The first four numeric columns are standard mIoU under compact and large official vocabularies. The final columns are diagnostic readouts, not replacement ranking metrics.}
\label{tab:e2-ambiguity}
\setlength{\tabcolsep}{2.5pt}
\begin{adjustbox}{max width=\textwidth}
\begin{tabular}{l c c c c c c p{0.22\textwidth}}
\toprule
Method & Ctx-59 mIoU\up & Ctx-459 mIoU\up & ADE-150 mIoU\up & ADE-847 mIoU\up & $\Delta_{\mathrm{vocab}}$\down & MCMR@0.5\down & Top mismatch pairs \\
\midrule
MaskCLIP & & & & & & & \\
CLIP-DIY & & & & & & & \\
SCLIP & & & & & & & \\
NACLIP & & & & & & & \\
ResCLIP & & & & & & & \\
ProxyCLIP & & & & & & & \\
CorrCLIP & & & & & & & \\
Trident & & & & & & & \\
OVDiff / FreeDA & & & & & & & \\
GroundingDINO + SAM2 & & & & & & & \\
SAM-AMG + SigLIP & & & & & & & \\
DINOv2 + SAM + SigLIP & & & & & & & \\
OVSeg / SAN / ODISE (ref.) & & & & & & & \\
\bottomrule
\end{tabular}
\end{adjustbox}
\end{table}

\subsection{E3: Cross-Dataset Generalization}

\begin{table}[H]
\centering
\scriptsize
\caption{E3 cross-dataset generalization. E3 uses the same standard metric as E1 and E2: target mIoU. Scenario difficulty comes from source-to-target transfer and dataset variation, not a different quality metric.}
\label{tab:e3-generalization}
\setlength{\tabcolsep}{2.5pt}
\begin{adjustbox}{max width=\textwidth}
\begin{tabular}{l p{0.22\textwidth} c c c c c c}
\toprule
Method & Source / training data & VOC-20 mIoU\up & Context-59 mIoU\up & ADE-150 mIoU\up & COCO-Stuff mIoU\up & Avg.\up & Worst\up \\
\midrule
\multicolumn{8}{l}{\textbf{Strict training-free rows}} \\
MaskCLIP & None & & & & & & \\
CLIP-DIY & None & & & & & & \\
SCLIP & None & & & & & & \\
NACLIP & None & & & & & & \\
ResCLIP & None & & & & & & \\
ProxyCLIP & None & & & & & & \\
CorrCLIP & None & & & & & & \\
Trident & None & & & & & & \\
FreeDA & None; offline diffusion prototypes only & & & & & & \\
GroundingDINO + SAM2 & None & & & & & & \\
SAM-AMG + SigLIP & None & & & & & & \\
DINOv2 + SAM + SigLIP & None & & & & & & \\
\midrule
\multicolumn{8}{l}{\textbf{Compact training-based references}} \\
OVSeg & Official released weights / COCO-style source & & & & & & \\
SAN & Official released weights / COCO-Stuff source & & & & & & \\
ODISE & Official released weights / COCO source & & & & & & \\
\bottomrule
\end{tabular}
\end{adjustbox}
\end{table}

\paragraph{E3 reporting rule.}
For the main E3 table, every method is evaluated with official target vocabularies and target mIoU. Hard-domain camouflaged-object transfer is still useful, but it is reported in the appendix because those targets require mask-only metrics and would break the unified E1--E3 metric structure.

\subsection{E4: Efficiency Comparison}

\begin{table}[H]
\centering
\scriptsize
\caption{E4 efficiency comparison. Exact numbers are filled after running every method on the same GPU, input size, and adapter settings. Training-based methods are intentionally limited to compact references.}
\label{tab:e4-efficiency}
\setlength{\tabcolsep}{3pt}
\begin{adjustbox}{max width=\textwidth}
\begin{tabular}{l c p{0.17\textwidth} c c p{0.36\textwidth}}
\toprule
Method & Training? & Train source & Train time & Test cost & Major modules / notes \\
\midrule
MaskCLIP & No & None & 0 & report & CLIP dense label extraction baseline. \\
CLIP-DIY & No & None & 0 & report & CLIP patch inference plus unsupervised localization prior. \\
SCLIP & No & None & 0 & report & CLIP with correlative self-attention. \\
NACLIP & No & None & 0 & report & CLIP neighbour-aware attention adaptation. \\
CLIPtrase & No & None & 0 & report & CLIP self-attention / dense inference modification. \\
SC-CLIP & No & None & 0 & report & Self-calibrated CLIP; no target training. \\
ResCLIP & No & None & 0 & report & Residual cross-correlation attention and semantic feedback refinement. \\
ProxyCLIP & No & None & 0 & report & CLIP plus VFM proxy attention. \\
CorrCLIP & No & None & 0 & report & CLIP with patch correlation reconstruction. \\
Trident & No & None & 0 & report & CLIP + DINO + SAM high-resolution training-free pipeline. \\
CASS & No & None & 0 & report & CLIP with VFM spectral object-context distillation at inference. \\
OVDiff & No & None & offline synth. & report & Diffusion-generated support/prototype stage; no target training. \\
FreeDA & No & None & offline proto. & report & Offline diffusion-augmented prototype generation. \\
GroundingDINO + SAM2 & No & None & 0 & report & Open-set detector plus SAM2 mask generation. \\
SAM-AMG + SigLIP & No & None & 0 & report & Class-agnostic SAM proposals plus SigLIP naming. \\
DINOv2 + SAM + SigLIP & No & None & 0 & report & DINOv2-guided proposals plus SigLIP naming. \\
OVSeg & Yes & COCO-style source & report & report & Mask-adapted CLIP trained reference. \\
SAN & Yes & COCO-Stuff source & report & report & Side adapter trained reference. \\
ODISE & Yes & COCO source & report & report & Diffusion-based open-vocabulary panoptic trained reference. \\
\bottomrule
\end{tabular}
\end{adjustbox}
\end{table}

\subsection{Diagnostic Probes}

\begin{table}[H]
\centering
\scriptsize
\caption{Failure-source diagnostic probes. These are not model ablations and are not separate benchmark tasks. They explain whether errors come from naming, localization, proposal recall, or mask--category consistency.}
\label{tab:diagnostic-probes}
\begin{adjustbox}{max width=\textwidth}
\begin{tabular}{p{0.22\textwidth} p{0.40\textwidth} p{0.18\textwidth} p{0.16\textwidth}}
\toprule
Diagnostic probe & Setting & Reported quantities & Ranking use \\
\midrule
GT-region naming diagnostic & Classify ground-truth regions with CLIP/SigLIP; localization is removed. & Top-1 / Top-5 naming accuracy & Diagnostic only \\
GT-text localization diagnostic & Segment or localize using the ground-truth class prompt; recognition is removed. & IoU / BIoU / mask metrics & Diagnostic only \\
Proposal-recall diagnostic & Choose the proposal with highest GT IoU from SAM-AMG or DINOv2+SAM; naming is removed. & Oracle IoU / BIoU & Diagnostic only \\
MCC consistency diagnostic & Compare the same proposal+naming pipeline with and without post-hoc mask--category consistency re-ranking. & Pair-consistency gain / failure cases & Diagnostic only \\
\bottomrule
\end{tabular}
\end{adjustbox}
\end{table}

%==============================================================
\clearpage

\section*{A Appendix Tables}

\begin{table}[H]
\centering
\scriptsize
\caption{Appendix hard-domain transfer stress test. These targets are not used to define the unified E1--E3 main metric because they require object-mask metrics rather than standard semantic mIoU.}
\label{tab:appendix-hard-domain}
\begin{adjustbox}{max width=\textwidth}
\begin{tabular}{l c c c c c c c}
\toprule
Method & OVCamo cIoU\up & OVCamo mask\up & CAMO mask\up & COD10K-Camo mask\up & NC4K mask\up & Hard avg\up & Notes \\
\midrule
MaskCLIP & & & & & & & \\
SCLIP & & & & & & & \\
NACLIP & & & & & & & \\
ProxyCLIP & & & & & & & \\
Trident & & & & & & & \\
FreeDA & & & & & & & \\
GroundingDINO + SAM2 & & & & & & & \\
SAM-AMG + SigLIP & & & & & & & \\
DINOv2 + SAM + SigLIP & & & & & & & \\
OVSeg / SAN / ODISE (ref.) & & & & & & & \\
\bottomrule
\end{tabular}
\end{adjustbox}
\end{table}

\paragraph{Hard-domain mask score.}
The hard-domain mask columns are filled with the standard camouflaged-object metric suite: IoU, $S_m$, $F^\omega_\beta$, $E_m$, and MAE. The compact table may show one representative mask score or an average after all numbers are available, while the spreadsheet keeps the full suite. OVCamo cIoU is reported only when the class-aware OVCamo vocabulary is used.

\begin{table}[H]
\centering
\scriptsize
\caption{E2 qualitative mismatch summary supporting MCMR@0.5. These entries are explanatory evidence, not additional benchmark metrics.}
\label{tab:appendix-ambiguity-graph}
\begin{adjustbox}{max width=\textwidth}
\begin{tabular}{l c c c p{0.34\textwidth}}
\toprule
Method group & Localized pairs & Mismatch pairs & MCMR@0.5 & Example mismatch pairs / communities \\
\midrule
Dense-map TF methods & & & & \\
Detector+SAM TF methods & & & & \\
Proposal+naming TF methods & & & & \\
MCC consistency diagnostics & & & & \\
Trained references & & & & \\
\bottomrule
\end{tabular}
\end{adjustbox}
\end{table}

\begin{table}[H]
\centering
\scriptsize
\caption{Spreadsheet-recorded exploratory readouts. Each method-dataset run produces all of these so they are available for offline analysis. They are promoted to the main paper only when they pass the promotion rule (cross-dataset stability, family-level discrimination, mIoU-orthogonal signal).}
\label{tab:appendix-supporting-readouts}
\begin{adjustbox}{max width=\textwidth}
\begin{tabular}{p{0.23\textwidth} p{0.44\textwidth} p{0.25\textwidth}}
\toprule
Readout & What to record & Possible use \\
\midrule
\textbf{E1 effectiveness diagnostics} & & \\
Per-class behavior & Per-class IoU, thing/stuff mIoU, pixel accuracy, mean class accuracy. & Explain which classes or category types drive mIoU changes. \\
Frequency-stratified mIoU & Head/mid/tail mIoU using training-frequency tertiles. & Detect head-class-only methods; long-tail behavior. \\
Object-scale stratified mIoU & Small/medium/large mIoU using COCO-style scale bins. & Reveal resolution / proposal-granularity issues. \\
Confidence calibration & ECE, Brier score, reliability diagrams. & Test whether TF methods are systematically miscalibrated. \\
Coverage / over-seg. rates & Empty-prediction rate, over-segmentation rate, prediction coverage. & Catch silent failures and proposal flooding. \\
\midrule
\textbf{E2 vocabulary / ambiguity diagnostics} & & \\
Vocabulary ambiguity (core) & MCMR@0.5, MCMR@0.75, top mismatch pairs, confusion communities. & Core E2 analysis. \\
Semantic-distance-weighted mismatch & Per-error weight by CLIP / WordNet distance between $\hat{y}_i$ and $y_i$. & Distinguish near-miss vs.\ wild-miss errors. \\
Hierarchical IoU & Partial credit when $\hat{y}_i$ is a WordNet hyper/hyponym of $y_i$. & Test whether plain mIoU underestimates OV semantic ability. \\
Synonym-collapsed mIoU & mIoU after merging WordNet-synset synonyms. & Quantify synonym competition vs.\ real ambiguity. \\
Prompt-sensitivity variance & mIoU std.\ across a fixed prompt-template bank. & Separate prompt-fragile from prompt-stable methods. \\
Vocabulary-size scaling curves & mIoU at \{59,150,300,459\} on Context, \{150,300,500,847\} on ADE. & Curve shape itself is a finding. \\
Region-level Top-$k$ naming & Top-1/Top-3/Top-5 accuracy on regions with $\mathrm{IoU}\ge 0.5$. & Separate "does not know" from "knows but ranks wrong." \\
GT oracle diagnostics & GT-region naming Top-1/Top-5; GT-text localization IoU / BIoU. & Separate naming and localization failures. \\
Proposal quality & BestIoU, Recall@0.5, Recall@0.7, number of proposals. & Ceiling for proposal-based pipelines. \\
MCC consistency & Before/after mIoU, before/after MCMR@0.5, qualitative examples. & Check whether post-hoc consistency reduces mismatch. \\
\midrule
\textbf{E3 generalization diagnostics} & & \\
Rank correlations & Kendall's $\tau$, Spearman's $\rho$ between E1 and E3 rankings; cross-target rank stability. & Hard evidence of generalization fragility. \\
Class-overlap stratified mIoU & Target mIoU split into source-overlapping vs.\ source-novel classes. & Expose trained-reference dependence on source vocabulary. \\
Domain-shift sensitivity & mIoU drop vs.\ DINO / CLIP source-to-target embedding distance. & Identify clean training-free signatures. \\
Worst-class within target & Minimum per-class IoU within each target. & Catch one-class collapses hidden by averages. \\
Per-target $\Delta$ MCMR & MCMR@0.5 transferred from E1/E2 datasets to E3 targets. & Separate localization vs.\ naming transfer failures. \\
\midrule
\textbf{E4 efficiency exploratory} & & \\
Efficiency details & Offline time, per-image inference time, memory, model calls, input resolution. & Build cost-performance analysis. \\
Normalized efficiency & mIoU per GFLOP, mIoU per second. & Single number for Pareto-frontier plots. \\
Offline amortization curve & Per-image total cost as a function of $N$ (eval-set size). & Fair treatment of diffusion-prototype methods. \\
Per-stage breakdown & Memory and latency per encoder/proposal/naming/post-proc.\ stage. & Identify the bottleneck stage per method. \\
\bottomrule
\end{tabular}
\end{adjustbox}
\end{table}

\begin{table}[H]
\centering
\scriptsize
\caption{Per-method exclusion or demotion grounds for rows outside the strict TF leaderboard.}
\label{tab:appendix-exclusion}
\begin{adjustbox}{max width=\textwidth}
\begin{tabular}{l c c c c p{0.34\textwidth}}
\toprule
Method & Task training? & Prompt/adapt. training? & MLLM call? & API-VLM? & Main reason \\
\midrule
OVSeg & Yes & Yes & No & No & Kept as compact trained reference only. \\
SAN & Yes & Yes & No & No & Kept as compact trained reference only. \\
ODISE & Yes & Yes & No & No & Kept as compact trained reference only. \\
OVCoser & Yes & Yes & No & No & OVCOS-specific; moved to hard-domain related work / appendix. \\
SuCLIP & Yes & Yes & No & No & OVCOS-specific; moved to hard-domain related work / appendix. \\
COCUS / cascaded VLM OVCOS & Yes & Yes & No & No & Related hard-domain work; not part of general OVOS main rows. \\
DSS / Discover-Segment-Select & No & No & Yes & No & Uses inference-time MLLM mask selection; related work only. \\
GenSAM / ProMaC & No & No & Yes & No & Inference-time generative MLLM; related work only. \\
API-VLM + SAM & No & No & Depends & Yes & Closed-source API; not reproducible under strict protocol. \\
\bottomrule
\end{tabular}
\end{adjustbox}
\end{table}

\clearpage
\begin{thebibliography}{99}
\small
\bibitem{revisitovs2025} Qiming Huang, Han Hu, and Jianbo Jiao. Revisit the Open Nature of Open Vocabulary Semantic Segmentation. ICLR, 2025.
\bibitem{ovparts2023} Meng Wei, Xiaoyu Yue, Wenwei Zhang, Shu Kong, Xihui Liu, and Jiangmiao Pang. OV-PARTS: Towards Open-Vocabulary Part Segmentation. NeurIPS Datasets and Benchmarks, 2023.
\bibitem{ovdg2026} Dong Zhao, Qi Zang, Nan Pu, Wenjing Li, Nicu Sebe, and Zhun Zhong. Open-Vocabulary Domain Generalization in Urban-Scene Segmentation. arXiv, 2026.
\bibitem{pascalvoc2010} Mark Everingham et al. The PASCAL Visual Object Classes (VOC) Challenge. IJCV, 2010.
\bibitem{pascalcontext2014} Roozbeh Mottaghi et al. The Role of Context for Object Detection and Semantic Segmentation in the Wild. CVPR, 2014.
\bibitem{ade20k2017} Bolei Zhou et al. Scene Parsing through ADE20K Dataset. CVPR, 2017.
\bibitem{cocostuff2018} Holger Caesar, Jasper Uijlings, and Vittorio Ferrari. COCO-Stuff: Thing and Stuff Classes in Context. CVPR, 2018.
\bibitem{clip2021} Alec Radford et al. Learning Transferable Visual Models From Natural Language Supervision. ICML, 2021.
\bibitem{siglip2023} Xiaohua Zhai et al. Sigmoid Loss for Language Image Pre-Training. ICCV, 2023.
\bibitem{dinov22023} Maxime Oquab et al. DINOv2: Learning Robust Visual Features without Supervision. TMLR, 2024.
\bibitem{sam2023} Alexander Kirillov et al. Segment Anything. ICCV, 2023.
\bibitem{sam22024} Nikhila Ravi et al. SAM 2: Segment Anything in Images and Videos. arXiv, 2024.
\bibitem{groundingdino2023} Shilong Liu et al. Grounding DINO: Marrying DINO with Grounded Pre-Training for Open-Set Object Detection. ECCV, 2024.
\bibitem{maskclip2022} Chong Zhou, Chen Change Loy, and Bo Dai. Extract Free Dense Labels from CLIP. ECCV, 2022.
\bibitem{clipdiy2024} Monika Wysoczanska, Michael Ramamonjisoa, Tomasz Trzcinski, and Oriane Simeoni. CLIP-DIY: CLIP Dense Inference Yields Open-Vocabulary Semantic Segmentation For-Free. WACV, 2024.
\bibitem{sclip2024} Feng Wang, Jieru Mei, and Alan Yuille. SCLIP: Rethinking Self-Attention for Dense Vision-Language Inference. ECCV, 2024.
\bibitem{naclip2025} Sina Hajimiri, Ismail Ben Ayed, and Jose Dolz. Pay Attention to Your Neighbours: Training-Free Open-Vocabulary Semantic Segmentation. WACV, 2025.
\bibitem{cliptrase2024} Tong Shao, Zhuotao Tian, Hang Zhao, and Jingyong Su. Explore the Potential of CLIP for Training-Free Open Vocabulary Semantic Segmentation. ECCV, 2024.
\bibitem{scclip2025} Sule Bai, Yong Liu, Yifei Han, Haoji Zhang, and Yansong Tang. Self-Calibrated CLIP for Training-Free Open-Vocabulary Segmentation. arXiv, 2024/2025.
\bibitem{resclip2025} Yiwen Yang et al. ResCLIP: Residual Attention for Training-free Dense Vision-language Inference. CVPR, 2025.
\bibitem{proxyclip2024} Mengcheng Lan et al. ProxyCLIP: Proxy Attention Improves CLIP for Open-Vocabulary Segmentation. ECCV, 2024.
\bibitem{corrclip2025} Dengke Zhang, Fagui Liu, and Quan Tang. CorrCLIP: Reconstructing Patch Correlations in CLIP for Open-Vocabulary Semantic Segmentation. ICCV, 2025.
\bibitem{trident2024} Yuheng Shi, Minjing Dong, and Chang Xu. Harnessing Vision Foundation Models for High-Performance, Training-Free Open Vocabulary Segmentation. arXiv, 2024.
\bibitem{cass2025} Chanyoung Kim, Dayun Ju, Woojung Han, Minchul Yang, and Seong Jae Hwang. Distilling Spectral Graph for Object-Context Aware Open-Vocabulary Semantic Segmentation. CVPR, 2025.
\bibitem{ovdiff2024} Laurynas Karazija, Iro Laina, Andrea Vedaldi, and Christian Rupprecht. Diffusion Models for Open-Vocabulary Segmentation. ECCV, 2024.
\bibitem{freeda2024} Luca Barsellotti, Roberto Amoroso, Marcella Cornia, Lorenzo Baraldi, and Rita Cucchiara. Training-Free Open-Vocabulary Segmentation with Offline Diffusion-Augmented Prototype Generation. CVPR, 2024.
\bibitem{ovseg2023} Feng Liang et al. Open-Vocabulary Semantic Segmentation with Mask-Adapted CLIP. CVPR, 2023.
\bibitem{san2023} Mengde Xu, Zheng Zhang, Fangyun Wei, Han Hu, Xiang Bai. Side Adapter Network for Open-Vocabulary Semantic Segmentation. CVPR, 2023.
\bibitem{odise2023} Jiarui Xu et al. Open-Vocabulary Panoptic Segmentation with Text-to-Image Diffusion Models. CVPR, 2023.
\bibitem{ovcoser2024} Youwei Pang, Xiaoqi Zhao, Jiaming Zuo, Lihe Zhang, and Huchuan Lu. Open-Vocabulary Camouflaged Object Segmentation. ECCV, 2024.
\bibitem{suclip2025} Ren et al. Seeing the Unseen: A Semantic Alignment and Context-Aware Prompt Framework for Open-Vocabulary Camouflaged Object Segmentation. ICCV, 2025.
\bibitem{camo2019} Trung-Nghia Le, Tam V. Nguyen, Zhongliang Nie, Minh-Triet Tran, and Akihiro Sugimoto. Anabranch Network for Camouflaged Object Segmentation. CVIU, 2019.
\bibitem{cod10k2020} Deng-Ping Fan et al. Camouflaged Object Detection. CVPR, 2020.
\bibitem{nc4k2021} Yunqiu Lv et al. Simultaneously Localize, Segment and Rank the Camouflaged Objects. CVPR, 2021.
\end{thebibliography}

\end{document}
